
\subsubsection{2.1 Adversarial Robustness Benchmarks}

Ribeiro et al. (2020) introduced CheckList, a behavioral testing paradigm evaluating NLP models across a matrix of linguistic capabilities and perturbation types. CheckList shifts the evaluation frame from aggregate accuracy to structured failure-mode analysis, but does not compare architecture families. Wang et al. (2021) released AdvGLUE, a multi-task adversarial benchmark derived from five GLUE tasks using 14 attack methods spanning character-level noise, word-level substitution, and semantic perturbations; it is the primary evaluation suite in this study. Nie et al. (2020) developed ANLI (Adversarial NLI), where examples are collected via iterative human-model adversarial interaction, providing a benchmark resistant to statistical artifacts from automatic perturbation methods. ANLI Round 3 (ANLI-R3, 1,200 examples) is used here as the human-crafted, surrogate-free evaluation partition.

A persistent limitation in prior use of these benchmarks is their application for leaderboard-style scalar comparison. The present work instead treats AdvGLUE and ANLI-R3 as sources of per-category $\Delta$* observations and subjects the resulting attack-category $\times$ model matrix to multivariate analysis.

\subsubsection{2.2 Architecture Comparison Studies}

Neerudu et al. (2023) compare BERT, GPT-2, and T5 on GLUE with 8 text perturbation types and find that GPT-2 representations are more robust than BERT and T5 across multiple perturbation categories. This directional finding is consistent with the $\eta^{2}=0.293$ result reported here and with the attention-topology hypothesis. However, Neerudu et al. evaluate only three models (one per family), report no effect size, apply no multivariate test, and do not use AdvGLUE or ANLI. The present work extends that directional observation with formal effect-size quantification, a nine-model pool, a six-category $\Delta$*-vector representation, and permutation MANOVA with statistical controls.

Zhang et al. (2026) find that decoder-only models show lower explanation flip rates than encoder baselines in enterprise NLP settings, consistent with architecture-family robustness differences in deployment contexts. Yang et al. (2024) compare LLM robustness under white-box attacks and observe that model size and structure affect robustness, without formal between-family effect-size quantification.

\subsubsection{2.3 LLM Trustworthiness Evaluation}

Sun et al. (2024, TrustLLM) evaluate 16 LLMs across six trustworthiness dimensions using more than 30 datasets. The robustness dimension is relevant to the present study, but TrustLLM does not stratify by architecture family or compute per-attack-category $\Delta$\textit{-vectors. Otmakhova et al. (2025, FLUKE) demonstrate that a model's ability to use a linguistic feature does not imply robustness to perturbations of that feature, motivation consistent with the present approach of measuring $\Delta$} independently of clean accuracy. Cohen (1988) provides the multivariate effect-size conventions and power analysis framework applied throughout this paper.

\textbf{Summary.} Prior work establishes directional architecture-family robustness differences and documents them across diverse benchmarks. The gap is a formal, effect-size-quantified, scale-matched, multivariate analysis of between-family $\Delta$*-vector clustering under statistical controls --- and a validated reusable pipeline enabling community replication.

